%=========================================================================
%  THEME 1 — "KOMPAKT"
%
%  Einspaltig mit Datums-Schiene links. Entspricht dem tabellarischen
%  Lebenslauf, wie ihn HR im DACH-Raum liest, in moderner Typografie:
%  Inter, Haarlinien, eine Akzentfarbe, konsequenter Flatterrand.
%  Lineare Lesereihenfolge — von Menschen und Parsern gleich gelesen.
%=========================================================================

\usepackage[left=1.7cm, right=1.7cm, top=1.25cm, bottom=1.4cm,
            headheight=14pt, headsep=12pt, footskip=1cm]{geometry}

%=========================================================================
%  GEMEINSAME BASIS — Pakete, Icons, Hilfsmakros.
%  Enthält bewusst KEINE Gestaltung. Alles Gestalterische (Farben, Schrift,
%  Raster, Abstände) liegt im jeweiligen Theme unter templates/themes/.
%=========================================================================

\usepackage[british, ngerman]{babel}
\usepackage{etoolbox}
\usepackage{xcolor}
\usepackage{fontspec}
\usepackage[final]{microtype}

%-------------------------------------------------------------------------
%  Echte Leerzeichen im PDF. TeX setzt Wortzwischenräume sonst
%  nur als Abstand. Ältere PDF-Leser (z.B. pdf.js 2.x im PDF-Viewer von
%  VS Code) und manche Bewerbungsportale lesen den Text dann ohne Leerzeichen
%  aus ("Bei500TaktenproMinute"). Mit 2 schreibt XeTeX die Leerzeichen als
%  Glyphen mit. Seitenzahlen und Blocksatz bleiben gleich, einzelne Absätze
%  brechen etwas anders um (weniger Trennungen).
%  Gemessen mit pdf.js 2.6, Anteil einzeln lesbarer Wörter: Briefe vorher
%  3-5 %, nachher 97-99 %. Im Lebenslauf stieg der Anteil ebenfalls deutlich.
%  Rest: Zeilen in der Kompetenzentabelle und einzelne Zeilen mit Bindestrich
%  verlieren die Zeichen weiterhin (XeTeX-intern, nicht weiter eingegrenzt).
%  pdftotext, pdf.js 4 und gängige Parser setzen dort ohnehin Leerzeichen aus
%  den Abständen, ohne Doppelungen.
%  ACHTUNG, XeTeX 0.999993 stürzt damit ab (Segmentation fault), wenn
%    - irgendwo \spaceskip ungleich null ist (deshalb der Umbau bei ragged2e
%      unten) oder
%    - \RaggedRight/\bodyflow mitten in einem schon begonnenen Absatz
%      eingeschaltet wird ("Text {\bodyflow ...}"). Immer am Absatzanfang.
%  Und: Farbwechsel mitten im Fliesstext in derselben Schrift verrutschen
%  (das farbige Wort bleibt schwarz, der Text danach wird farbig). Abhilfe:
%  \kern0pt an beiden Enden des farbigen Stücks, z.B.
%  {\color{red}\kern0pt Text\kern0pt} (so \platz in feedback.tex.
%  Farbe zusammen mit Schriftwechsel (Ortsangaben, Symbole) ist nicht betroffen.
%  Stürzt ein Build nach einer Vorlagenänderung ab, zuerst diese beiden Punkte
%  prüfen. Notausgang: hier \XeTeXinterwordspaceshaping=0 setzen.
%-------------------------------------------------------------------------
\XeTeXinterwordspaceshaping=2
% Als Zeichen geschrieben wird nur ein normaler Wortabstand. Nach Grossbuch-
% staben (sfcode 999) und Satzzeichen (ohne \frenchspacing) setzt TeX einen
% leicht abweichenden Abstand und XeTeX lässt das Zeichen weg ("SAPHANACloud").
% Deshalb alle auf 1000 -- im Deutschen ohnehin üblich, im Satzbild unsichtbar.
\count255=`A
\loop\sfcode\count255=1000 \ifnum\count255<`Z \advance\count255 by 1 \repeat
\sfcode`Ä=1000 \sfcode`Ö=1000 \sfcode`Ü=1000
\AtBeginDocument{\frenchspacing}
% Unsichtbares Leerzeichen ohne Breite: dort, wo im Satz kein Wortabstand
% steht und Leser trotzdem einen brauchen -- zwischen zwei Spalten
% ("EntwicklungPython"), nach Aufzählungszeichen und am Schriftwechsel nach
% den Leads ("Prüfstandsdatenplattform:Von"). Ändert das Layout nicht.
\newcommand{\cvrealspace}{\ifnum\XeTeXcharglyph"20>0
    \rlap{\XeTeXglyph\XeTeXcharglyph"20\relax}\fi}
\usepackage{tikz}
\usetikzlibrary{calc}
\usepackage{graphicx}
\usepackage{enumitem}
\usepackage{fancyhdr}
\usepackage{tabularx}
\usepackage{array}
\usepackage[fixed]{fontawesome5}
\usepackage{paracol}
\usepackage{setspace}
\usepackage{ragged2e}
% ragged2e hält im Flattersatz die Wortabstände starr, indem es bei jedem
% Schriftwechsel \spaceskip auf den natürlichen Wortabstand setzt. Zusammen
% mit echten Leerzeichen (\XeTeXinterwordspaceshaping=2, oben) stürzt XeTeX
% 0.999993 dabei ab. Gleicher Effekt ohne \spaceskip:
% Dehnung und Stauchung des Wortabstands in der Schrift selbst auf null setzen
% (\fontdimen3/4). Der Wortabstand ist damit derselbe wie bei ragged2e.
% Nebenwirkung: Die Änderung gilt für die ganze Satzlauf-Instanz der Schrift.
% Unschädlich, weil die Briefe im Blocksatz vor dem Lebenslauf gesetzt werden
% bzw. in eigenen Läufen (make dossier).
\makeatletter
\renewcommand{\@raggedtwoe@everyselectfont}{%
    \spaceskip\z@
    \if@raggedtwoe@spaceskip
        \fontdimen3\font=\z@
        \fontdimen4\font=\z@
    \fi}
\makeatother
\usepackage{needspace}
\usepackage{accsupp}

% Textschicht für Bewerbungsportale. SmartRecruiters & Co.
% lesen den Lebenslauf aus dem PDF aus. Ohne diese Zeile kamen gesperrte
% Überschriften als "BER UFS E R FAH R UNG" und jeder Bindestrich als
% geschützter Bindestrich (U+2011) an. \XeTeXgenerateactualtext legt zu
% jedem Textstück den Originaltext ab; am Druckbild ändert sich nichts.
\ifdefined\XeTeXgenerateactualtext\XeTeXgenerateactualtext=1\fi


%-------------------------------------------------------------------------
%  Sprache
%
%  \cvlanguage ist ein babel-Sprachname und wird VOR dem Theme gesetzt --
%  in applications/<app>/meta.tex, denn die Sprache gehoert zur Ziel-
%  bewerbung und nicht zur Gestaltung. Voreinstellung ngerman: alle
%  bestehenden Bewerbungen bauen unveraendert, ohne dass ihre meta.tex
%  etwas davon wissen muss. Die babel-Optionen oben laden alle drei
%  Sprachen; die letzte ist die Hauptsprache, \selectlanguage schaltet um.
%  Fuer Englisch "british" verwenden, nicht "english": babel warnt vor der
%  mehrdeutigen Option, und british liefert das europaeische Datumsformat.
%-------------------------------------------------------------------------

\providecommand{\cvlanguage}{ngerman}

%-------------------------------------------------------------------------
%  Beschriftungen der Abschnitte -- Voreinstellung Deutsch.
%  Eine uebersetzte Inhaltsdatei (templates/themes/cv-content-en.tex)
%  setzt sie per \renewcommand um; kein Theme enthaelt mehr festen Text.
%-------------------------------------------------------------------------

%  Briefdatum. Voreinstellung ist babels \today. Die englische Inhalts-
%  datei setzt es um: babels british schreibt "16th September 2026", und
%  die Ordnungszahl ist im Geschaeftsbrief unueblich.
\providecommand{\lettertoday}{\today}

\providecommand{\LblSkills}{Kompetenzen}
\providecommand{\LblExperience}{Berufserfahrung}
\providecommand{\LblEducation}{Ausbildung}
\providecommand{\LblFurther}{Weiterbildung}
\providecommand{\LblLanguages}{Sprachen}
\providecommand{\LblContact}{Kontakt}
\providecommand{\LblProfile}{Profil}
\providecommand{\LblCV}{Lebenslauf}
\providecommand{\LblPage}{Seite}

\setlength{\parindent}{0pt}
\setlength{\parskip}{0pt}
\emergencystretch=2em

%-------------------------------------------------------------------------
%  Typografische Hygiene
%-------------------------------------------------------------------------

% Trennung in Überschriften/Titeln unterdrücken — verhindert Unfälle wie
% "MUSTERMANN TECH-|NIK AG" quer über den Zeilenumbruch.
\newcommand{\nohyph}{\hyphenpenalty=10000\exhyphenpenalty=10000\hbadness=10000}

% Eigennamen, die nie getrennt werden sollen. \hyphenation traegt die
% Ausnahmen nur in die gerade aktive Sprache ein -- deshalb als Makro und
% nach \selectlanguage noch einmal aufgerufen.
\newcommand{\cvhyphenexceptions}{%
    % Eigene Firmen-, Orts- und Produktnamen hier ergänzen. Ohne Bindestrich
    % = nie trennen, mit Bindestrich = nur an dieser Fuge trennen.
    \hyphenation{ANSYS FastAPI Team-center Work-bench}}
\cvhyphenexceptions
% \frenchspacing NACH \selectlanguage: babel setzt den Satzabstand je
% Sprache, und british stellt den weiten englischen Abstand nach dem Punkt
% ein. Der faellt im Brief als Loecher auf und passt nicht zu den deutschen
% Dokumenten, die ngerman ohnehin auf frenchspacing stellt -- fuer Deutsch
% ist die Zeile daher wirkungslos.
\AtBeginDocument{\expandafter\selectlanguage\expandafter{\cvlanguage}%
    \cvhyphenexceptions\frenchspacing}

% Flatterrand statt Blocksatz: in schmalen Spalten erzeugt Blocksatz
% Trennungskaskaden und Löcher. \RaggedRight erlaubt Trennung noch,
% hält den Grauwert aber ruhig. \cvbodyhyphenpenalty kann die Inhaltsdatei
% überschreiben: 10000 schaltet die Trennung ab  --
% getrennte Wörter kamen in Bewerbungsportalen als "Haushaltsvoll- zugs-
% programm" an, weil die Portale den Text aus dem PDF auslesen.
\providecommand{\cvbodyhyphenpenalty}{200}
% Ohne Trennung dehnt \raggedright statt \RaggedRight: ragged2e lässt eine
% Zeile vor einem langen Wort lieber überlaufen ("Prozessdokumentation"
% ragte 13,6 pt in den Rand).
\newcommand{\bodyflow}{\RaggedRight\hyphenpenalty=\cvbodyhyphenpenalty\relax
    \ifnum\cvbodyhyphenpenalty<10000 \else\raggedright\fi}

%-------------------------------------------------------------------------
%  Icons
%-------------------------------------------------------------------------

% Die Symbolschrift hat keine Unicode-Zuordnung; ausgelesen wurde jedes
% Symbol zu einem Störzeichen vor der Adresse oder Telefonnummer ("Ȱ", "Л").
% \icotext gibt ihnen einen leeren Ersatztext -- Portale sehen nur die Angabe.
\newcommand{\icotext}[1]{\BeginAccSupp{method=plain,ActualText={}}#1\EndAccSupp{}}
\newcommand{\icoloc}{\icotext{\faMapMarker*}}
\newcommand{\icophone}{\icotext{\faPhone*}}
\newcommand{\icomail}{\icotext{\faEnvelope}}
\newcommand{\icolinked}{\icotext{\faLinkedin}}
\newcommand{\icocal}{\icotext{\faBirthdayCake}}
\newcommand{\icoheart}{\icotext{\faHeart}}
\newcommand{\icocar}{\icotext{\faCar}}

%-------------------------------------------------------------------------
%  Sechster Kontaktplatz — z.B. LinkedIn, GitHub oder der Führerschein.
%  Die Inhaltsdatei setzt Symbol, Text und optional eine URL; ist die URL
%  leer, steht der Text ohne Verlinkung da.
%-------------------------------------------------------------------------
\providecommand{\CVlinkicon}{\icolinked}
\providecommand{\CVlinktext}{}
\providecommand{\CVlinkurl}{}
\newcommand{\cvlink}{\ifdefempty{\CVlinkurl}{\CVlinktext}{\href{\CVlinkurl}{\CVlinktext}}}

%-------------------------------------------------------------------------
%  Haken für Abschnitte, die nur eine Person braucht (z.B. Referenzen).
%  Wird am Ende des Lebenslaufs ausgeführt; die Inhaltsdatei füllt ihn.
%-------------------------------------------------------------------------
\providecommand{\cvextrasections}{}

%-------------------------------------------------------------------------
%  Master-Lebenslauf und optionale Bausteine.
%
%  Die Inhaltsdatei hält ALLE Angaben. Was im Dossier erscheint, steuert
%  die meta.tex der Bewerbung:
%      \cvshow{schluessel}   zeigt einen Baustein, der sonst wegbleibt
%      \cvhide{schluessel}   lässt einen Baustein weg, der sonst erscheint
%  Die Schlüssel stehen im Kopf von themes/cv-content.tex. application.tex
%  und preview.tex definieren \cvshow und \cvhide, BEVOR sie die meta.tex
%  lesen; hier stehen sie nur als Rückfallebene für andere Treiber.
%
%  In der Inhaltsdatei:
%      \cvopt{schluessel}{...}        nur mit \cvshow oder im Master
%      \cvstd{schluessel}{...}        immer, ausser mit \cvhide (im Master immer)
%      \cvifopt{schluessel}{dann}{sonst}
%  \def\cvmaster{} (make cv-master über templates/cv-master.tex) zeigt
%  alles: Das ist der Master-Lebenslauf, aus dem je Stelle reduziert wird.
%-------------------------------------------------------------------------
\providecommand{\cvshow}[1]{\expandafter\def\csname cvshow@#1\endcsname{}}
\providecommand{\cvhide}[1]{\expandafter\def\csname cvhide@#1\endcsname{}}
\newcommand{\cvifopt}[3]{\ifdef{\cvmaster}{#2}{\ifcsdef{cvshow@#1}{#2}{#3}}}
\newcommand{\cvopt}[2]{\cvifopt{#1}{#2}{}}
\newcommand{\cvstd}[2]{\ifdef{\cvmaster}{#2}{\ifcsdef{cvhide@#1}{}{#2}}}

%-------------------------------------------------------------------------
%  Achter Kontaktplatz: GitHub (Schlüssel github in der
%  Inhaltsdatei, gedacht für Software- und Datenstellen). Leer gelassen
%  erscheint er in keinem Theme.
%-------------------------------------------------------------------------
\providecommand{\CVgithubtext}{}
\providecommand{\CVgithuburl}{}
\newcommand{\icogithub}{\icotext{\faGithub}}
\newcommand{\cvgithub}{\ifdefempty{\CVgithuburl}{\CVgithubtext}{\href{\CVgithuburl}{\CVgithubtext}}}

%-------------------------------------------------------------------------
%  Haken für einen Abschnitt direkt nach der Berufserfahrung (z.B. "Open Source & Projekte"). Leer = kein Abschnitt. Jedes
%  Theme ruft ihn in \makecvbody auf.
%-------------------------------------------------------------------------
\providecommand{\cvprojectsection}{}

%-------------------------------------------------------------------------
%  Zwischenüberschrift innerhalb einer Station (z.B. "Entwicklung" und
%  "Versuch & Prüfung"). Steht zwischen
%  zwei cvbullets-Listen und bleibt mit dem ersten Punkt darunter zusammen.
%  Vorgabe für alle Themes; kompakt setzt eine eigene Fassung.
%-------------------------------------------------------------------------
\providecommand{\cvsubhead}[1]{%
    \par\Needspace*{5\baselineskip}\addvspace{4pt}%
    {\nohyph\mediumfont #1}\par\nopagebreak
}

%-------------------------------------------------------------------------
%  Siebter Kontaktplatz — Nationalität und Bewilligung. In der Schweiz
%  gehört die Angabe bei ausländischen Bewerbenden in die Personalien
%  (Aufenthaltsbewilligung B/C, für Liechtenstein die Grenzgänger-
%  bewilligung). Leer gelassen erscheint die Zeile nicht. 
%-------------------------------------------------------------------------
\providecommand{\CVnationality}{}
\newcommand{\icoflag}{\icotext{\faPassport}}

%-------------------------------------------------------------------------
%  Neunter Kontaktplatz — Bewilligung als eigener Platz:
%  "Deutsch, Niederlassungsbewilligung C" in einem Platz lief im Kompakt-
%  Theme bis in den GitHub-Platz daneben. Mit \CVpermit stehen Nationalität
%  und Bewilligung getrennt. Leer gelassen bleibt alles wie bisher.
%-------------------------------------------------------------------------
\providecommand{\CVpermit}{}
\newcommand{\icopermit}{\icotext{\faIdCard}}

%-------------------------------------------------------------------------
%  Schrift des Brieftexts. Voreinstellung: die \normalsize des Themes,
%  also dieselbe Grösse wie der Lebenslauf. Die Inhaltsdatei kann den
%  Brief grösser setzen — bei Briefen von 200 bis 300 Wörtern bleibt die
%  Seite in 9,4 pt sonst zur Hälfte leer . Wird am Ende
%  von \makeletterhead ausgeführt, nach dem \parskip des Themes.
%-------------------------------------------------------------------------
\providecommand{\letterbodyfont}{}

%-------------------------------------------------------------------------
%  Unterschrift als Bild. Leer = keine. Die Briefe rufen \lettersignature
%  zwischen Grussformel und Namen auf; ohne Bild bleibt nur der Abstand.
%-------------------------------------------------------------------------
\providecommand{\CVsignature}{}
\newcommand{\lettersignature}{%
    \ifdefempty{\CVsignature}{\vspace{0.3cm}}{%
        \par\vspace{2pt}\noindent\includegraphics[height=0.95cm]{\CVsignature}\par\vspace{1pt}}}

%-------------------------------------------------------------------------
%  Bild-Helfer
%-------------------------------------------------------------------------
% Quadratischer Ausschnitt mit abgerundeten Ecken.
%   #1 Kantenlänge des Rahmens   #2 Eckenradius
% Zwei Stellschrauben kommen aus der Inhaltsdatei, weil sie zum Foto
% gehören und nicht zum Theme — beide relativ zur Kantenlänge, damit
% derselbe Wert in allen Themes stimmt, egal wie gross der Rahmen ist:
%   \CVphotozoom   Bildbreite als Vielfaches der Kantenlänge (1 = bündig)
%   \CVphotoshift  Motivversatz; negativ schiebt das Bild nach unten,
%                  das Gesicht rutscht also von oben in die Mitte
\providecommand{\CVphotozoom}{1}
\providecommand{\CVphotoshift}{0}
%   \CVphotomirror 1 = wie aufgenommen, -1 = seitenverkehrt (damit die Person zum
%                  Text hin blickt, wenn das Foto rechts steht)
\providecommand{\CVphotomirror}{1}
\newcommand{\photobox}[2]{%
    \begin{tikzpicture}
        \pgfmathsetlengthmacro{\pbhalf}{0.5*#1}
        \pgfmathsetlengthmacro{\pbwide}{\CVphotozoom*#1}
        \pgfmathsetlengthmacro{\pbshift}{\CVphotoshift*#1}
        \clip[rounded corners=#2] (-\pbhalf,-\pbhalf) rectangle (\pbhalf,\pbhalf);
        \node[anchor=center, inner sep=0] at (0,\pbshift)
            {\scalebox{\CVphotomirror}[1]{\includegraphics[width=\pbwide]{\CVphoto}}};
    \end{tikzpicture}%
}

%-------------------------------------------------------------------------
%  Zweiter Zeitraum derselben Rolle beim selben Arbeitgeber.
%
%  Wer eine Funktion zweimal innehatte, bekommt EINEN Eintrag mit zwei
%  Zeiträumen — nicht zwei Einträge mit identischer Aufgabenliste. Der
%  jüngere Zeitraum steht oben und bestimmt, wo der Eintrag im Lebenslauf
%  einsortiert wird; der ältere folgt als zweite Zeile in der Datumsspalte.
%
%      \cvrole{Immobilien-Bewirtschafterin}{Ing.~Bau~AG}{Vaduz}%
%             {2022 -- 2023\cvalso{2018 -- 2019}}
%
%  Bewusst ohne "und": die Datumsspalte des Akzent-Themes ist 2,2 cm breit,
%  "und 2018 -- 2019" bricht dort auf drei Zeilen um. Dass es zwei Einsätze
%  sind, sagt ohnehin die Zeile unter dem Rollentitel.
%-------------------------------------------------------------------------
\newcommand{\cvalso}[1]{\\ #1}

%-------------------------------------------------------------------------
%  Organisationszeile "Firma · Ort" — beide Teile optional, damit Einträge
%  ohne Arbeitgeber (Familienphase, Auszeit) keinen leeren Trenner zeigen.
%  Nutzt \orgfont, \metafont und die Farbe muted aus dem jeweiligen Theme.
%-------------------------------------------------------------------------
\newcommand{\cvorgline}[2]{%
    % Keine Silbentrennung: Firmen- und Ortsnamen gehören nicht zerlegt.
    % "Zwickau, Deutsch-|land" liest sich wie ein Setzfehler. Umbrüche an
    % den Wortzwischenräumen bleiben möglich.
    \nohyph
    \ifblank{#1}{%
        \ifblank{#2}{}{{\metafont\color{muted}#2}}%
    }{%
        {\orgfont#1}%
        \ifblank{#2}{}{\,{\metafont\color{muted}\textbullet}\,{\metafont\color{muted}#2}}%
    }%
}

%-------------------------------------------------------------------------
%  Links — klickbar, ohne Rahmen
%-------------------------------------------------------------------------
\usepackage[hidelinks, bookmarks=false, unicode]{hyperref}
\urlstyle{same}

%-------------------------------------------------------------------------
%  Weitere optionale Angaben. Leer gelassen erscheinen sie nirgends.
%    \CVkeywords      Stichwortzeile unter dem Titel (kompakt, rueckgrat)
%    \CVwebtext/-url  eigene Website als weiterer Link (rueckgrat)
%    \CVphotoaspect   Höhe des Fotos im Verhältnis zur Breite (rueckgrat),
%                     1 = quadratisch, 1.25 = Hochformat 4:5
%-------------------------------------------------------------------------
\providecommand{\CVkeywords}{}
\providecommand{\CVwebtext}{}
\providecommand{\CVweburl}{}
\newcommand{\cvweb}{\ifdefempty{\CVweburl}{\CVwebtext}{\href{\CVweburl}{\CVwebtext}}}
\providecommand{\CVphotoaspect}{1}

\providecommand{\LblProjects}{Projekte}
\providecommand{\LblReferences}{Referenzen}
\providecommand{\LblEnclosures}{Beilagen}
\providecommand{\LblContents}{Inhalt}
\providecommand{\LblApplicationAs}{Bewerbung als}
\providecommand{\LblCoverLetter}{Anschreiben}
\providecommand{\LblMotivation}{Motivationsschreiben}

%-------------------------------------------------------------------------
%  Beilagenzeile am Ende des Anschreibens. Erscheint nur, wenn die meta.tex
%  \newcommand{\enclosures}{Lebenslauf, Arbeitszeugnisse, Diplome} setzt.
%-------------------------------------------------------------------------
\newcommand{\letterenclosures}{\ifdefined\enclosures
    \par\vspace{0.5cm}\noindent{\metafont\color{muted}\LblEnclosures: \enclosures}\par\fi}

%-------------------------------------------------------------------------
%  Deckblatt (\cvshow{coverpage} in der meta.tex): Briefkopf des Themes,
%  die Zielstelle in der Seitenmitte, unten das Inhaltsverzeichnis des
%  Dossiers. \coverpagecontents setzt application.tex je nach Schaltern.
%-------------------------------------------------------------------------
\providecommand{\coverpagecontents}{\LblCV}
% Grösse der Zielstelle auf dem Deckblatt, vom Theme überschreibbar
\providecommand{\coverpagetitlesize}{\fontsize{24pt}{29pt}\selectfont}
\newcommand{\makecoverpage}{%
    \makeletterhead
    \setlength{\parskip}{0pt}%
    \vspace*{\fill}%
    \noindent{\metafont\color{muted}\MakeUppercase{\LblApplicationAs}}\par\vspace{8pt}%
    \noindent{\nohyph\displayfont\coverpagetitlesize\color{headercolor}\positionname\par}%
    \vspace{6pt}%
    \noindent{\fontsize{13pt}{17pt}\selectfont\companyname\par}%
    \vspace*{\fill}\vspace*{\fill}%
    \noindent{\metafont\color{muted}\MakeUppercase{\LblContents}}\par\vspace{4pt}%
    \noindent{\color{rule}\rule{\linewidth}{0.5pt}}\par\vspace{6pt}%
    \noindent\coverpagecontents\par
    \newpage
}


%-------------------------------------------------------------------------
%  Farben
%-------------------------------------------------------------------------

\definecolor{ink}{HTML}{1D2429}      % Fliesstext / Titel
\definecolor{accent}{HTML}{1F5D74}   % Petrol — Sektionen, Daten, Akzentlinien
\definecolor{muted}{HTML}{6B7680}    % Sekundärangaben
\definecolor{rule}{HTML}{DCE1E5}     % Haarlinien

% Kompatibilität zu den bestehenden Brief-Dateien
\colorlet{headercolor}{ink}
\colorlet{linecolor}{accent}

%-------------------------------------------------------------------------
%  Schrift — Inter durchgehend, Ziffern tabellarisch (Jahreszahlen fluchten)
%-------------------------------------------------------------------------

\setmainfont{Inter}[
    Extension      = .otf,
    UprightFont    = *-Regular,
    BoldFont       = *-SemiBold,
    ItalicFont     = *-Italic,
    BoldItalicFont = *-SemiBoldItalic,
    Numbers        = Monospaced,
    Ligatures      = TeX,
]
% \sffamily muss dieselbe Schrift liefern — die Briefdateien verwenden es
\setsansfont{Inter}[
    Extension = .otf, UprightFont = *-Regular, BoldFont = *-SemiBold,
    ItalicFont = *-Italic, Numbers = Monospaced, Ligatures = TeX]
\newfontfamily\displayfont{Inter}[
    Extension = .otf, UprightFont = *-SemiBold, BoldFont = *-Bold,
    Numbers = Monospaced, Ligatures = TeX]
\newfontfamily\mediumfont{Inter}[
    Extension = .otf, UprightFont = *-Medium, BoldFont = *-SemiBold,
    Numbers = Monospaced, Ligatures = TeX]

% Grundschriftgrad des Dokuments
\makeatletter
\renewcommand\normalsize{\@setfontsize\normalsize{9.2}{12.8}}
\makeatother
\AtBeginDocument{\normalsize\color{ink}}

\newcommand{\namefont}{\displayfont\fontsize{25pt}{25pt}\selectfont
    \addfontfeature{LetterSpace=-2.2}}
\newcommand{\rolelabelfont}{\mediumfont\fontsize{9.5pt}{12pt}\selectfont
    \addfontfeature{LetterSpace=16}}
\newcommand{\metafont}{\fontsize{8.3pt}{11pt}\selectfont}
\newcommand{\datefont}{\mediumfont\fontsize{8.5pt}{11pt}\selectfont}
\newcommand{\rolefont}{\displayfont\fontsize{10.8pt}{13pt}\selectfont}
\newcommand{\orgfont}{\mediumfont\fontsize{9.4pt}{12pt}\selectfont}
\newcommand{\notefont}{\fontsize{8.7pt}{11.5pt}\selectfont\itshape}
\newcommand{\lead}[1]{{\mediumfont#1\cvrealspace}}

% Grundlinien-Schutz: \color erzeugt in vertikalem Modus ein Whatsit; steht
% es als erstes im \vtop einer [t]-Minipage, verliert die Box ihre Höhe und
% die Zeile rutscht um die Oberlänge nach unten. \leavevmode verhindert das.
\newcommand{\vtopsafe}{\leavevmode}

% Zwischenspeicher für das Kurzprofil — siehe \makecvheader
\newsavebox{\summarybox}

%-------------------------------------------------------------------------
%  Sektionen
%-------------------------------------------------------------------------

% Optionales Argument: Platzbedarf in Zeilen, damit die Überschrift nicht allein
% am Seitenfuss steht, bei einem Lebenslauf mit drei Seiten: "Ausbildung" stand
% unten auf Seite 2, der Inhalt auf Seite 3). Voreinstellung 3 = Überschrift, Linie
% und eine Zeile (Kompetenzen, Sprachen, Referenzen). Folgt ein Stationseintrag
% (\cvrole verlangt selbst 5 Zeilen), verlangt der Aufrufer 8: \cvsection[8]{...}.
% Ein fester hoher Wert für alle Überschriften schob die Zweiseiter auf drei Seiten.
\newcommand{\cvsection}[2][3]{%
    \par\Needspace*{#1\baselineskip}\addvspace{10pt}%
    \noindent{\mediumfont\fontsize{9pt}{9pt}\selectfont
        \addfontfeature{LetterSpace=18}\color{accent}\MakeUppercase{#2}}%
    % \nointerlineskip: ohne ihn setzt TeX die Linie als eigene Zeile im
    % vollen Grundlinienabstand, und der Abstand wird deutlich grösser als
    % die 4pt hier. Versalien haben keine Unterlängen, der Abstand ist also exakt.
    \par\nointerlineskip\vspace{4pt}%
    \noindent{\color{rule}\rule{\linewidth}{0.5pt}}%
    \par\vspace{5pt}%
}

%-------------------------------------------------------------------------
%  Einträge — Datum in der Schiene, Inhalt rechts.
%  Als LaTeX-Liste realisiert, damit lange Einträge sauber über
%  Seitenumbrüche laufen; Minipages könnten das nicht.
%-------------------------------------------------------------------------

\newenvironment{cvexperience}{%
    \begin{list}{}{%
        \setlength{\leftmargin}{3.15cm}%
        \setlength{\labelwidth}{2.65cm}%
        \setlength{\labelsep}{0.5cm}%
        \setlength{\itemindent}{0pt}%
        \setlength{\listparindent}{0pt}%
        \setlength{\rightmargin}{0pt}%
        \setlength{\topsep}{0pt}%
        \setlength{\partopsep}{0pt}%
        \setlength{\parsep}{0pt}%
        \setlength{\itemsep}{8pt}%
        \renewcommand{\makelabel}[1]{%
            \parbox[t]{\labelwidth}{\raggedright\vtopsafe
                \datefont\color{accent}\strut##1}\hfill}%
    }%
}{\end{list}\par\addvspace{1pt}}

\newcommand{\cvrole}[4]{%
    \Needspace*{5\baselineskip}% Station nicht als Waise an den Seitenfuss
    \item[#4]%
    {\nohyph\rolefont\strut#1}\par\vspace{2pt}%
    \cvorgline{#2}{#3}%
    \par\vspace{2.5pt}%
}

\newcommand{\cvnote}[1]{\vspace{-1pt}{\bodyflow\notefont\color{muted}#1\par}\vspace{1pt}}

\newcommand{\cvplain}[3]{%
    \item[#2]\bodyflow
    {\nohyph\orgfont\strut#1}\,{\metafont\color{muted}\textbullet}\,#3\par%
}

% Ein Punkt wird nie über einen Seitenumbruch geteilt (\interlinepenalty);
% umbrochen wird nur zwischen zwei Punkten. beginpenalty hält die Liste an
% ihrer Überschrift bzw. an der Firmenzeile der Station.
\newenvironment{cvbullets}{%
    \begin{itemize}[beginpenalty=10000,
                    leftmargin=13pt, labelwidth=7pt, labelsep=6pt, labelindent=0pt,
                    itemindent=0pt, topsep=2.5pt, itemsep=2.5pt,
                    parsep=0pt, partopsep=0pt, align=left,
                    label={\color{accent}\textendash\cvrealspace}]
        \bodyflow\interlinepenalty=10000\relax
    }{\end{itemize}}

% Zwischenüberschrift in einer Station: Firmenschrift in Akzentfarbe, damit sie
% sich von der Firmenzeile (Grundfarbe) und den Leads (fett) abhebt.
\renewcommand{\cvsubhead}[1]{%
    \par\Needspace*{5\baselineskip}\addvspace{5pt}%
    {\nohyph\orgfont\color{accent}#1}\par\nopagebreak\vspace{0.5pt}%
}

%-------------------------------------------------------------------------
%  Fakten-Zeilen (Kompetenzen, Sprachen) — dieselbe Schiene wie die
%  Berufserfahrung, damit das Dokument ein einziges Raster hat.
%-------------------------------------------------------------------------

\newcommand{\cvfact}[2]{%
    \noindent
    \begin{minipage}[t]{2.65cm}\raggedright\vtopsafe
        \datefont\color{accent}\strut#1\cvrealspace\end{minipage}%
    \hspace{0.5cm}%
    \begin{minipage}[t]{\dimexpr\linewidth-3.15cm\relax}\bodyflow\strut#2\end{minipage}%
    \par\vspace{4pt}%
}

\newcommand{\cvtag}[1]{#1\quad}  % Chips braucht dieses Theme nicht

%-------------------------------------------------------------------------
%  Kopf des Lebenslaufs
%-------------------------------------------------------------------------

\newcommand{\cbox}[2]{%
    \begin{minipage}[t]{0.333\linewidth}\vtopsafe\metafont\color{muted}
        \strut#1~~#2\end{minipage}%
}

\newcommand{\contactgrid}{%
    \noindent
    \cbox{\icoloc}{\CVaddress}%
    \cbox{\icophone}{\CVphone}%
    \cbox{\icomail}{\href{mailto:\CVmail}{\CVmail}}\par\vspace{3pt}
    \noindent
    \cbox{\CVlinkicon}{\cvlink}%
    \cbox{\icocal}{\CVbirth}%
    \cbox{\icoheart}{\CVcivil}\par
    % Dritte Zeile: Nationalität, Bewilligung und, bei Software-/
    % Datenstellen, GitHub. Alle optional, die Zeile entfällt,
    % wenn alle leer sind.
    \ifboolexpr{test {\ifdefempty{\CVnationality}} and test {\ifdefempty{\CVpermit}}
                and test {\ifdefempty{\CVgithubtext}}}{}{%
        \vspace{3pt}\noindent
        \ifdefempty{\CVnationality}{}{\cbox{\icoflag}{\CVnationality}}%
        \ifdefempty{\CVpermit}{}{\cbox{\icopermit}{\CVpermit}}%
        \ifdefempty{\CVgithubtext}{}{\cbox{\icogithub}{\cvgithub}}\par}%
}

% Kopf mit zweispaltigen Kontaktangaben: Sie stehen unter dem Namen, das Foto läuft rechts
% über die ganze Kopfhöhe. Eingeschaltet über \cvheadertwocol, das die Inhaltsdateien setzen.
% \newcommand{\cvheaderclassic}{} in der meta.tex schaltet für eine Bewerbung zurück auf den
% früheren Kopf (drei Spalten über die ganze Breite, Foto 3.4 cm).
% \photowide ist die Kantenlänge des Fotos: 4.5 cm passt zu vier Kontaktzeilen, bei drei
% Zeilen sind es 3.9 cm.
\newlength{\photowide}\setlength{\photowide}{4.5cm}
\newcommand{\cboxtwo}[2]{%
    \begin{minipage}[t]{0.5\linewidth}\vtopsafe\metafont\color{muted}
        \strut#1~~#2\end{minipage}%
}
\newcommand{\contactgridtwo}{%
    \noindent
    \cboxtwo{\icoloc}{\CVaddress}%
    \cboxtwo{\icophone}{\CVphone}\par\vspace{3pt}
    \noindent
    \cboxtwo{\icomail}{\href{mailto:\CVmail}{\CVmail}}%
    \cboxtwo{\CVlinkicon}{\cvlink}\par\vspace{3pt}
    \noindent
    \cboxtwo{\icocal}{\CVbirth}%
    \cboxtwo{\icoheart}{\CVcivil}\par
    \ifboolexpr{test {\ifdefempty{\CVnationality}} and test {\ifdefempty{\CVpermit}}}{}{%
        \vspace{3pt}\noindent
        \ifdefempty{\CVnationality}{}{\cboxtwo{\icoflag}{\CVnationality}}%
        \ifdefempty{\CVpermit}{}{\cboxtwo{\icopermit}{\CVpermit}}\par}%
    \ifdefempty{\CVgithubtext}{}{\vspace{3pt}\noindent\cboxtwo{\icogithub}{\cvgithub}\par}%
}

\newcommand{\makecvheader}{%
  \ifdefined\cvheadertwocol
    \noindent
    \begin{minipage}[b]{\dimexpr\linewidth-\photowide-0.7cm\relax}
        {\namefont\nohyph\color{ink}\CVname}\par\vspace{6pt}
        {\rolelabelfont\color{accent}\MakeUppercase{\CVrole}}\par
        \ifdefempty{\CVkeywords}{}{\vspace{3pt}{\metafont\color{muted}\CVkeywords}\par}%
        \vspace{9pt}
        \noindent{\color{rule}\rule{\linewidth}{0.5pt}}\par\vspace{6pt}
        \contactgridtwo
    \end{minipage}%
    \hfill
    \begin{minipage}[b]{\photowide}
        \raggedleft\cvstd{photo}{\photobox{\photowide}{7pt}}
    \end{minipage}
    \par\vspace{11pt}
  \else
    \noindent
    \begin{minipage}[b]{\dimexpr\linewidth-3.9cm\relax}
        {\namefont\nohyph\color{ink}\CVname}\par\vspace{6pt}
        {\rolelabelfont\color{accent}\MakeUppercase{\CVrole}}
    \end{minipage}%
    \hfill
    \begin{minipage}[b]{3.4cm}
        \raggedleft\cvstd{photo}{\photobox{3.4cm}{6pt}}
    \end{minipage}
    \par\vspace{9pt}
    \noindent{\color{rule}\rule{\linewidth}{0.5pt}}\par\vspace{6pt}
    \contactgrid
    \par\vspace{9pt}
  \fi
    % Kurzprofil, links geklammert von einer Akzentlinie. Der Text wird
    % zuerst in eine Box gesetzt, dann bekommt die Linie deren gemessene
    % Höhe: eine fest gesetzte Länge würde je nach Textmenge über den
    % Absatz hinausragen oder zu kurz bleiben.
    \sbox{\summarybox}{%
        \begin{minipage}[t]{\dimexpr\linewidth-20pt\relax}
            \vtopsafe\bodyflow\fontsize{9.3pt}{13pt}\selectfont
            \strut\ContentSummary\strut
        \end{minipage}}%
    \noindent
    {\color{accent}\rule[-\dp\summarybox]{1.6pt}%
        {\dimexpr\ht\summarybox+\dp\summarybox\relax}}%
    \hspace{9pt}\usebox{\summarybox}\par
    % Groesser als die 10pt, die \cvsection selbst setzt: der Profilkasten
    % endet auf einem \strut und damit ohne Unterlaenge, wodurch der gleiche
    % Zahlenwert hier optisch enger wirkt als zwischen zwei Sektionen.
    % \addvspace nimmt das Maximum, deshalb muss der Wert oberhalb von 10pt
    % liegen, um ueberhaupt zu greifen.
    \vspace{17pt}
}

%-------------------------------------------------------------------------
%  Seitenstil — dezente Kopfzeile ab Seite 2
%-------------------------------------------------------------------------

\fancypagestyle{cvstyle}{%
    \fancyhf{}
    \renewcommand{\headrulewidth}{0pt}
    \renewcommand{\footrulewidth}{0pt}
    \fancyhead[L]{\metafont\color{muted}\CVname~~\textbullet~~\LblCV}
    \fancyhead[R]{\metafont\color{muted}\LblPage\ \thepage}
    \fancyhead[C]{}
}
\fancypagestyle{cvfirst}{\fancyhf{}\renewcommand{\headrulewidth}{0pt}}

%-------------------------------------------------------------------------
%  Briefkopf (Anschreiben / Motivationsschreiben)
%-------------------------------------------------------------------------

\newcommand{\makeletterhead}{%
    \pagestyle{cvfirst}\thispagestyle{cvfirst}%
    \setlength{\parskip}{7.5pt}% Brieftext: Absätze durch Weissraum getrennt
    \noindent
    \begin{minipage}[b]{\dimexpr\linewidth-5.4cm\relax}
        {\displayfont\fontsize{19pt}{19pt}\selectfont
         \addfontfeature{LetterSpace=-1.5}\color{ink}\CVname}\par\vspace{5pt}
        {\mediumfont\fontsize{8.8pt}{11pt}\selectfont
         \addfontfeature{LetterSpace=14}\color{accent}\MakeUppercase{\CVrole}}
    \end{minipage}%
    \hfill
    \begin{minipage}[b]{5.2cm}
        \raggedleft\metafont\color{muted}\setlength{\parskip}{0pt}
        \CVaddress\par\CVphone\par\href{mailto:\CVmail}{\CVmail}
    \end{minipage}
    \par\vspace{8pt}
    \noindent{\color{accent}\rule{\linewidth}{1pt}}
    \par\vspace{12pt}%
    \letterbodyfont
}

\newcommand{\lettersection}[1]{%
    \par\addvspace{11pt}%
    \noindent{\mediumfont\fontsize{9.5pt}{9.5pt}\selectfont
        \addfontfeature{LetterSpace=14}\color{accent}\MakeUppercase{#1}}%
    \par\vspace{4pt}\noindent{\color{rule}\rule{\linewidth}{0.5pt}}\par\vspace{6pt}%
}

%-------------------------------------------------------------------------
%  Rumpf des Lebenslaufs
%-------------------------------------------------------------------------

\newcommand{\makecvbody}{%
    \pagestyle{cvstyle}\thispagestyle{cvfirst}
    \makecvheader

    \cvsection{\LblSkills}
    \ContentSkills

    \cvsection[8]{\LblExperience}
    \ContentExperience

    % Open Source & Projekte (nur, wenn die Inhaltsdatei den Haken füllt)
    \cvprojectsection

    \cvsection[8]{\LblEducation}
    \ContentEducation

    % Leere Weiterbildung druckt keine Überschrift
    \ifdefempty{\ContentFurther}{}{\cvsection[8]{\LblFurther}\ContentFurther}

    \cvsection{\LblLanguages}
    \ContentLanguages

    \cvextrasections
}
